\documentclass[10pt]{article}
\usepackage[T1]{fontenc}
\usepackage[utf8]{inputenc}
\usepackage[french]{babel}
\usepackage{amsmath,amssymb,mathrsfs,mathpazo}
\usepackage{geometry}
\geometry{a4paper,margin=2cm}
\usepackage{graphicx,array,tabularx,tabularray,colortbl,multirow,multicol,booktabs,enumitem,pifont,ulem,fancybox,diagbox,cancel}
\usepackage{pstricks,pstricks-add,pst-plot,pst-tree,pst-node,pst-text,pst-func,pst-math,pst-bspline,pst-3dplot,pst-eucl,pst-solides3d}
\let\pscalculate\undefined
\usepackage{pst-bezier}
\usepackage{tikz,pgfplots,tkz-tab}
\usetikzlibrary{arrows,arrows.meta,patterns,calc,positioning,trees,decorations.markings,shapes.geometric,shapes.misc,shapes.arrows,intersections,decorations.pathmorphing,decorations.pathreplacing,backgrounds,shadows,babel}
\pgfplotsset{compat=1.18}
\usepgfplotslibrary{fillbetween}
\usepackage[np]{numprint}
\usepackage{listings,makecell,pas-tableur,twoopt}
\let\barre\undefined
\usepackage[pstricks]{tablvar}
\definecolor{darkgreen}{RGB}{0,100,0}
\providecommand{\red}{\color{red}}
\providecommand{\blue}{\color{blue}}
\providecommand{\green}{\color{green}}
\providecommand{\black}{\color{black}}
\providecommand{\white}{\color{white}}
\providecommand{\gray}{\color{gray}}
\providecommand{\cyan}{\color{cyan}}
\providecommand{\magenta}{\color{magenta}}
\providecommand{\yellow}{\color{yellow}}
\providecommand{\violet}{\color{violet}}
\providecommand{\orange}{\color{orange}}
\providecommand{\euro}{\texteuro}
\providecommand{\EUR}{\texteuro}
\providecommand{\ieme}{\textsuperscript{e}}
\providecommand{\degres}{\ensuremath{^\circ}}
\providecommand{\degree}{\ensuremath{^\circ}}
\providecommand{\vv}[1]{\overrightarrow{#1}}
\providecommand{\Coord}[1]{#1}
\providecommand{\decosix}{}
\providecommand{\starredbullet}{\ensuremath{\bullet}}
\pagestyle{empty}
\setlength{\parindent}{0pt}
\usepackage{adjustbox}\DeclareUnicodeCharacter{2236}{:}\DeclareUnicodeCharacter{2190}{\ensuremath{\leftarrow}}
\providecommand{\R}{\mathbb{R}}
\providecommand{\N}{\mathbb{N}}
\providecommand{\Z}{\mathbb{Z}}
\providecommand{\Q}{\mathbb{Q}}
\providecommand{\D}{\mathbb{D}}
\providecommand{\C}{\mathbb{C}}
\providecommand{\e}{{\mathrm{e}}}
\providecommand{\vect}[1]{\overrightarrow{#1}}
\providecommand{\vectt}[1]{\overrightarrow{\mathrm{#1}}}
\providecommand{\vvt}[1]{\overrightarrow{\mathrm{#1}}}
\providecommand{\np}[2][]{{#2~#1}}
\providecommand{\vv}[1]{{\overrightarrow{#1}}}
\renewcommand{\d}{{\mathrm{d}}}
\providecommand{\degres}{{^\circ}}
\providecommand{\degre}{{^\circ}}
\providecommand{\blue}{\color{blue}}
\providecommand{\red}{\color{red}}
\providecommand{\enskip}{\quad}
\providecommand{\llbracket}{\lbrack\!\lbrack}
\providecommand{\rrbracket}{\rbrack\!\rbrack}
\providecommand{\euro}{{\text{€}}}
\providecommand{\no}{{\text{n°}}}
\providecommand{\interval}[2]{{[#1;#2]}}
\providecommand{\intervFO}[2]{{[#1;#2[}}
\providecommand{\dotfill}{\ldots}

\providecommand{\ds}{\displaystyle}
\providecommand{\textding}[1]{#1}
\providecommand{\Oij}{\(\mathrm{O};\vec{\imath},\vec{\jmath}\)}
\providecommand{\Oijk}{\(\mathrm{O};\vec{\imath},\vec{\jmath},\vec{k}\)}
\providecommand{\Ouv}{\(\mathrm{O};\vec{u},\vec{v}\)}
\providecommand{\intervFF}[2]{\ensuremath{[#1;#2]}}
\renewcommand{\euro}{\ifmmode\text{\texteuro}\else\texteuro\fi}
\DeclareUnicodeCharacter{00B0}{\ifmmode^{\circ}\else\textdegree\fi}
\DeclareUnicodeCharacter{20AC}{\ifmmode\text{\texteuro}\else\texteuro\fi}
% Énoncés : Éducation nationale. Transcription et corrigés : APMEP, Denis Vergès et contributeurs. Découpage et classement : CoopMaths, Sébastien Lozano, Éric Elter et contributeurs. Adaptation de la présentation : Maths Explorer.
% CC BY-SA : https://coopmaths.fr/www/about/
% https://coopmaths.fr/alea/static/e3c/2020/tex/e3c_2020_00_sujet30_2_cor.tex
\begin{document}
\noindent\textit{APMEP / CoopMaths. CC BY-SA. Présentation adaptée par Maths Explorer.}\par\medskip
\subsection*{1.}

En dérivant le produit, toutes les fonctions étant dérivables :
\[
f'(x) = 2 e^x + (2x - 1) e^x = e^x(2 + 2x - 1) = (2x + 1) e^x.
\]

\subsection*{2.}

On sait que, quel que soit \(x \in \mathbb{R}\), \(e^x > 0\) : le signe de \(f'(x)\) est celui de \(2x + 1\) qui s'annule pour :
\[
2x + 1 = 0 \quad \Longleftrightarrow \quad x = -\dfrac{1}{2}.
\]
On en déduit :
\begin{itemize}
    \item \(f'(x) > 0\) sur l'intervalle \(\left]-\dfrac{1}{2}\,;\,+\infty\right[\) ;
    \item \(f'(x) < 0\) sur l'intervalle \(\left]-\infty\,;\,-\dfrac{1}{2}\right[\).
\end{itemize}

\subsection*{3.}

La question précédente montre que la fonction \(f\) est décroissante sur l'intervalle \(\left]-\infty\,;\,-\dfrac{1}{2}\right[\), et croissante sur l'intervalle \(\left]-\dfrac{1}{2}\,;\,+\infty\right[\), avec un minimum en \(x = -\dfrac{1}{2}\), où :
\[
f\left(-\dfrac{1}{2}\right) = \left(2 \times \left(-\dfrac{1}{2}\right) - 1\right) e^{-\frac{1}{2}} = -2 e^{-\frac{1}{2}} \approx -1,213.
\]
D'où le tableau de variations sur \(\mathbb{R}\) :

\begin{center}
\begin{tikzpicture}
\tkzTabInit[lgt=3.5, espcl=3]{$x$ / 1, {Signe de $f'(x)$} / 1, {$f$} / 2}{${-\infty}$, ${-\dfrac{1}{2}}$, ${+\infty}$}
\tkzTabLine{,-,0,+,}
\tkzTabVar{+/{$$},-/{$\approx -1,213$},+/{$$}}{/}
\end{tikzpicture}
\end{center}

\subsection*{4.}

\(C\) coupe l'axe des ordonnées en des points dont l'ordonnée est nulle, donc tels que :
\[
(2x - 1) e^x = 0 \quad \Longleftrightarrow \quad 2x - 1 = 0 \quad (\text{car } e^x \neq 0) \text{ d'où } x = \dfrac{1}{2}.
\]

\(C\) coupe l'axe des ordonnées au point \(\left(\dfrac{1}{2}\,;\,0\right)\).

\subsection*{5.}

\(M(x\,;\,y) \in T \Longleftrightarrow y - f(0) = f'(0)(x - 0).\)

Avec \(f(0) = -e^0 = -1\) et \(f'(0) = 1e^0 = 1\), on obtient :

\(M(x\,;\,y) \in T \Longleftrightarrow y - (-1) = 1(x - 0) \Longleftrightarrow y = x - 1.\)

\end{document}
